\chapter{Lecture 2: Redirection, Globbing, Pipes and egrep}
\begin{center}
\textit{January 9, 2018}
\end{center}

\section{cat and Keyboard Signals}
\code{cat} stands for \emph{concatenate}; with file arguments it displays the
contents of the files one after another. With no arguments it echoes stdin to
stdout.

\begin{itemize}
  \item \textbf{Ctrl+C} --- kill the running program.
  \item \textbf{Ctrl+D} --- send the end-of-file (EOF) signal.
\end{itemize}

\section{Input and Output Redirection}

\begin{definition}
\textbf{Output redirection} sends the output produced by a program to a file:
\code{command arg1 arg2 ... > filename}.

\textbf{Input redirection} sends the contents of a file to a program as its
input: \code{command arg1 arg2 ... < filename}.
\end{definition}

\subsection*{\code{cat sample.txt} vs.\ \code{cat < sample.txt}}
The output is identical, but the mechanism is different:
\begin{itemize}
  \item \code{cat sample.txt}: \code{cat} is given a filename, \emph{cat} opens
        the file and displays its contents.
  \item \code{cat < sample.txt}: the \emph{shell} opens the file and sends its
        contents to \code{cat}'s stdin; \code{cat} never sees a filename.
\end{itemize}
\begin{center}
\begin{tikzpicture}[>=Stealth,font=\small,
  prog/.style={draw,fill=gray!15,minimum width=1.6cm,minimum height=0.9cm},
  file/.style={draw,fill=yellow!10,minimum height=0.7cm}]
  \node[prog] (c1) at (0,0) {\code{cat}};
  \node[file] (f1) at (-4.2,0) {\code{sample.txt}};
  \draw[->] (c1.west) -- node[above,font=\scriptsize]{opens (argument)} (f1.east);
  \node[font=\scriptsize,anchor=west] at (1.2,0) {\code{cat sample.txt}};
  \node[prog] (c2) at (0,-1.6) {\code{cat}};
  \node[file] (f2) at (-4.2,-1.6) {\code{sample.txt}};
  \node[draw,dashed,rounded corners,font=\scriptsize] (sh) at (-4.2,-2.6) {shell opens it};
  \draw[->] (f2.east) -- node[above,font=\scriptsize]{stdin} (c2.west);
  \draw[dashed] (sh) -- (f2);
  \node[font=\scriptsize,anchor=west] at (1.2,-1.6) {\code{cat < sample.txt}};
\end{tikzpicture}
\end{center}
The difference is visible with \code{wc}:
\begin{lstlisting}[style=bash]
$ wc sample.txt
 10  52 300 sample.txt
$ wc < sample.txt
 10  52 300
\end{lstlisting}

Input and output redirection can be combined (each used once):
\begin{lstlisting}[style=bash]
cmd arg1 arg2 ... < inputfile > outputfile
cat -n < in.txt > out.txt    # like cp, but numbers the lines
\end{lstlisting}

\section{Behind the Scenes: Three Streams}
Every Linux process (\code{ls}, \code{cat}, a C/C++ program, \dots) has three
streams associated with it.

\begin{center}
\begin{tabular}{@{}p{0.29\textwidth}p{0.29\textwidth}p{0.29\textwidth}@{}}
\toprule
stdin (standard input) & stdout (standard output) & stderr (standard error) \\
\midrule
Default: keyboard & Default: screen & Default: screen \\
Change with \code{<} & Change with \code{>} & Change with \code{2>} \\
 & Buffered & Not buffered \\
\bottomrule
\end{tabular}
\end{center}

\begin{aside}
\textbf{Buffering:} writing in blocks is faster than writing one character at a
time (you will see why in CS 350). stderr is unbuffered so error messages show up
immediately.
\end{aside}

\begin{lstlisting}[style=bash]
cmd >> out.txt               # append to the end of out.txt
cmd < in 1> out 2>&1         # stdout to out, stderr to the same place as stdout
\end{lstlisting}
\code{\&1} means ``wherever stream 1 (stdout) is currently going''.

\section{Wildcards and Globbing}
\begin{lstlisting}[style=bash]
ls *.txt
\end{lstlisting}
\code{*.txt} is a \textbf{globbing pattern}: \code{*} matches any sequence of
characters, so this lists all files whose names end with \code{.txt}.

The \emph{shell} intercepts the pattern:
\begin{enumerate}
  \item it finds the files that match the pattern, and
  \item replaces the pattern by those file names.
\end{enumerate}
So \code{ls} actually receives \code{ls sample.txt sample2.txt}. Because the
shell does the work, globbing works with any command:
\begin{center}
\begin{tikzpicture}[>=Stealth,font=\small,node distance=0.8cm]
  \node[draw,fill=yellow!10] (t) {\code{ls *.txt}};
  \node[draw,rounded corners,fill=blue!6,right=of t,align=center] (s) {shell\\ \scriptsize expands the glob};
  \node[draw,fill=yellow!10,right=of s] (t2) {\code{ls sample.txt sample2.txt}};
  \node[draw,fill=gray!15,right=of t2] (l) {\code{ls}};
  \draw[->] (t) -- (s); \draw[->] (s) -- (t2); \draw[->] (t2) -- (l);
\end{tikzpicture}
\end{center}
\begin{lstlisting}[style=bash]
cat *.txt
echo *.txt      # a cheap way to see what a pattern expands to
rm *.txt
rm -rf *        # careful!!
\end{lstlisting}
Other glob characters: \code{?} matches exactly one character, and \code{[abc]}
matches one of the listed characters.

\begin{important}[Quoting]
Single quotes and double quotes suppress globbing: \code{echo "*.txt"} prints
\code{*.txt}.
\end{important}

\section{Pipes}
We can send the output (stdout) of one program to the input (stdin) of another
with a \textbf{pipe}: \code{cmd1 | cmd2}.

\begin{example}
How many words occur in the first 20 lines of \code{sample.txt}?
\begin{lstlisting}[style=bash]
head -20 sample.txt | wc -w
\end{lstlisting}
\end{example}

\begin{example}
Files \code{words*.txt} contain lists of words, one per line. Print a
duplicate-free list of all the words in all the files.
\begin{lstlisting}[style=bash]
cat words*.txt | sort | uniq
\end{lstlisting}
\begin{center}
\begin{tikzpicture}[>=Stealth,font=\small,node distance=1.4cm,
  prog/.style={draw,fill=gray!15,minimum width=1.4cm,minimum height=0.8cm}]
  \node[prog] (a) {\code{cat}};
  \node[prog,right=of a] (b) {\code{sort}};
  \node[prog,right=of b] (c) {\code{uniq}};
  \node[right=of c] (d) {screen};
  \node[draw,fill=yellow!10,left=0.9cm of a,align=center,font=\scriptsize] (f) {\code{words1.txt}\\ \code{words2.txt}};
  \draw[->] (f) -- (a);
  \draw[->] (a) -- node[above,font=\scriptsize]{stdout} node[below,font=\scriptsize]{stdin} (b);
  \draw[->] (b) -- node[above,font=\scriptsize]{stdout} node[below,font=\scriptsize]{stdin} (c);
  \draw[->] (c) -- node[above,font=\scriptsize]{stdout} (d);
\end{tikzpicture}
\end{center}
\code{uniq} only removes \emph{adjacent} duplicates, which is why we
\code{sort} first.
\end{example}

\section{Embedded Commands}
Use \code{\$(cmd)} to embed the output of a command inside another command:
\begin{lstlisting}[style=bash]
echo "Today is $(date) and I am $(whoami)"
\end{lstlisting}
\begin{itemize}
  \item Double quotes allow embedded commands (and variables) but suppress
        globbing; they also keep the result as a \emph{single} argument.
  \item Single quotes suppress embedded commands as well:
        \code{echo 'Today is \$(date)'} prints the text literally.
  \item Without quotes, each whitespace-separated word is a separate argument.
\end{itemize}

\section{egrep}
\code{egrep pattern file(s)} prints the lines that contain a substring matching
the given pattern (a regular expression).
\begin{lstlisting}[style=bash]
egrep cs246 index.shtml | wc -l           # number of lines mentioning cs246
egrep "cs246|CS246" index.shtml | wc -l   # either spelling
\end{lstlisting}
\code{|} in a regular expression means ``or''. Use double quotes around patterns
containing characters that are special to the shell (\code{|}, \code{*},
\code{?}, \dots), otherwise the shell interprets them first.

\section*{Summary}
\begin{itemize}
  \item \code{>}, \code{>>}, \code{<}, \code{2>}, \code{2>\&1} redirect streams;
        the shell does the redirection, not the program.
  \item stdout is buffered, stderr is not.
  \item Globbing (\code{*}, \code{?}, \code{[...]}) is expanded by the shell;
        quotes suppress it.
  \item Pipes chain programs; \code{\$(...)} embeds a command's output.
  \item \code{egrep} searches for lines matching a regular expression.
\end{itemize}
